% Part: second-order-logic
% Chapter: metatheory
% Section: introduction

\documentclass[../../../include/open-logic-section]{subfiles}

\begin{document}

\olfileid{sol}{met}{int}

\olsection{تعارف}

پہلے درجے دی منطق دیاں کئی چنگیاں خاصیتاں نیں۔ اسیں جاندے آں کہ ایہ
فیصلہ کرن جوگ نہیں، پر گھٹ توں گھٹ ایہ مسلّماتی طور تے متعین کیتی جا سکدی
اے۔ یعنی پہلے درجے دی منطق لئی صحت مند تے تمام ثبوتی نظام موجود نیں۔ اوہناں توں اک !!a{derivability} تعلق~$\Proves$ ملدا اے۔ ایس دی خاصیت ایہ اے کہ !!{sentence}s دے کسے وی سیٹ~$\Gamma$ تے کسے وی
!!{sentence}~$!A$ لئی، $\Gamma \Entails !A$ اُدوں تے صرف اُدوں جدوں
$\Gamma \Proves !A$ ہووے۔ ایس توں خاص طور تے ایہ نکلدا اے کہ پہلے
درجے دی منطق دے منطقی طور تے معتبر جملے !!{computably
  enumerable} نیں۔ اک قابلِ حساب فنکشن~$f\colon \Nat \to
\Sent[L]$ موجود اے۔ ایس~$f$ دیاں قدراں عین~$\Lang{L}$ دے سارے تے صرف اوہی
!!{sentence}s نیں جیہڑے منطقی طور تے معتبر نیں۔ ایہ ایس لئی اے کہ
!!{derivation}s دی مؤثر گنتی وار فہرست بن سکدی اے۔ ایس فہرست وچ
جیہڑا اخذ اک~!!{sentence} !!{derive} کردا اے، اوہنوں اوسے~!!{sentence}
نال جوڑ دیندے آں۔ دوجے درجے دی منطق دی اظہاری طاقت پہلے درجے نالوں ودھ
اے، ایس لئی عام طور تے اوہدے سارے معتبر جملے لبھنے ودھ اوکھے نیں۔
اصل وچ اسیں دکھاواں گے کہ دوجے درجے دی منطق نہ صرف ناقابلِ فیصلہ اے
بلکہ ایس دے معتبر جملے !!{computably
  enumerable} وی نہیں۔ ایس دا مطلب اے کہ دوجے درجے دی منطق لئی کوئی
ایہو جیہا صحت مند تے تمام ثبوتی نظام نہیں ہو سکدا جیہدے اخذ مؤثر
طور تے فہرست بن سکدے ہون (بھاویں صحت مند پر نامکمل ثبوتی نظام موجود
نیں تے حقیقت وچ تحقیق دے اہم موضوع نیں)۔

پہلے درجے دی منطق دیاں ہور دو خاصیتاں وی نیں: ایہ محدود گواہی دی
خاصیت رکھدی اے (جے !!{sentence}s دے کسے سیٹ~$\Gamma$ دا ہر محدود
ذیلی سیٹ پورا ہون جوگا اے، تاں خود $\Gamma$ وی پورا ہون جوگا اے)،
تے ایس اُتے لووینہائم--سکولم دا قضیہ لاگو ہوندا اے (جے $\Gamma$
دا اک لامتناہی نمونہ ہووے، تاں اوہدا اک !!a{denumerable} نمونہ وی ہوندا اے)۔ ایہ دونیں نتیجے دوجے درجے دی
منطق لئی ناکام ہوندے نیں۔ وجہ فیر ایہی اے کہ دوجے درجے دی منطق
!!{domain}s دے آکار بارے اوہ حقائق بیان کر سکدی اے جیہڑے پہلے درجے
دی منطق نہیں کر سکدی۔

\end{document}
